% Source: ACM_Templates.pdf, PDF pp. 35-37.
% Corrections and validation: docs/validation.md.
\section{高斯消元：实数与异或方程组}
\label{algo:gauss}
\subsection*{原理：主元与自由变量}
对增广矩阵 $[A\mid b]$ 做交换两行、某行乘非零常数、某行加另一行倍数等初等行变换，解集保持不变。依次按列寻找主元，将该列其他行消成零，即可得到便于读解的形式。

\textbf{行指针和列指针必须分开。}某列在剩余行中找不到主元时，该列对应自由变量，跳过该列但不能消耗一行。消元结束后，设系数矩阵有 $m$ 列、秩为 $r$：
\begin{itemize}
    \item 存在系数全零而常数非零的一行：无解。
    \item 无矛盾且 $r=m$：唯一解。
    \item 无矛盾且 $r<m$：存在自由变量；实数域有无穷多解，自由度为 $m-r$。
\end{itemize}
把所有自由变量设为 $0$，再读取主元行，可得到一组特解。若要全部解，将每个自由变量依次设为 $1$、其他设为 $0$，对齐次系统求相应方向向量，再与特解作线性组合。

\subsection*{实数消元的使用教程}
\InlineCode{gauss(a,ans,eps)} 接收从 $0$ 开始的 $n\times(m+1)$ 增广矩阵，要求 $n\geq1$，末列为常数；矩阵按值传入，原数据保留。返回 $0$ 表示无解，$1$ 表示唯一解，$2$ 表示多解；有解时 \InlineCode{ans} 存一组特解，按未知数顺序从 $0$ 开始。

每次选取剩余行中绝对值最大的主元，以减轻数值误差；绝对值不超过 \InlineCode{eps} 时视为零。

时间复杂度为 $O(nm\min(n,m))$，额外空间为 $O(nm)$。固定绝对误差阈值适合量级适中的竞赛数据，数值量级跨度大或病态系统需另作缩放与误差分析。
\begin{lstlisting}[style=mtycpp,numbers=none]
vector<vector<double>> a={{1,1,3},{2,-1,0}};
vector<double> ans;
int status=gauss(a,ans);
\end{lstlisting}

\needspace{7\baselineskip}
\subsection*{模板代码}
\CodeFile{code/number-theory/gauss.cpp}

\subsection*{异或方程组}
每个变量为 $0$ 或 $1$，加法在 $\mathbb F_2$ 上进行，行加减都变成异或，不再使用浮点数。寻找值为 $1$ 的主元，随后把其他行该列的 $1$ 异或消去。

判定方式与上面相同，但有自由变量时解的数量是\textbf{有限的} $2^{m-r}$，不是无穷多。秩应按实际选中的主元计数，不能简单依赖消元后对角线上的 $1$。

\subsection*{异或模板的使用教程}
\InlineCode{xor_gauss(a,ans,rank)} 的矩阵与返回状态约定同上，元素必须为 $0$ 或 $1$；通过引用返回秩。若只需方案数，有解时根据自由变量数 $m-\InlineCode{rank}$ 计算 $2^{m-r}$，按题意决定取模或大整数。
\begin{lstlisting}[style=mtycpp,numbers=none]
vector<vector<int>> a={{1,1,0},{0,1,1}};
vector<int> ans;
int rank;
int status=xor_gauss(a,ans,rank);
\end{lstlisting}

\needspace{7\baselineskip}
\subsection*{模板代码}
\CodeFile{code/number-theory/xor_gauss.cpp}
